
\section{平方和分解与均值最优}
对任意实数 $c$，写 $y_i-c=(y_i-\bar y)+(\bar y-c)$，平方后相加。交叉项包含 $\sum_i(y_i-\bar y)=0$，因此
\[
\sum_i(y_i-c)^2=\sum_i(y_i-\bar y)^2+n(\bar y-c)^2.
\]
这一个恒等式同时解释常数回归基线、回归树叶均值和K均值更新。重要的学习方式是识别它在不同算法中的共同结构，而不把每章的均值当作互不相关的规定。

若有非负权重 $w_i$且总权重大于0，相应最优常数是 $\sum_iw_iy_i/\sum_iw_i$。权重可能表示频次、测量可靠性或损失代价，不同解释需要不同依据，不能仅为提高分数随意制造权重。

\section{链式法则的局部计算}
若 $z=wx+b$、$h=\phi(z)$、$q=vh+c$、$L=(q-y)^2/2$，则
\[
\frac{\partial L}{\partial w}=(q-y)\,v\,\phi'(z)\,x.
\]
四个因子分别来自损失对输出、输出对隐藏量、隐藏量对预激活、预激活对权重。若一个量影响多条路径，应将各路径贡献相加。梯度必须针对同一组参数计算；先改 $v$再用新 $v$算 $w$的梯度，就不是同一次前向状态的普通梯度更新。

中心差分可用来核对解析梯度，但要避开不可微点和极端步长。核对一个小模型比仅看大模型训练损失下降更容易定位符号或维度错误，因为许多错误程序也可能偶然使某个损失下降。

\section{相关性与平均的方差}
两个随机量的协方差为 $\operatorname{Cov}(A,B)=E[(A-EA)(B-EB)]$。在方差有限且非零时，相关系数是协方差除以两个标准差的乘积。它描述共同线性波动，相关性为0不一般等于独立。

对平均 $\bar Z=B^{-1}\sum_bZ_b$，展开平方后会出现所有交叉协方差。忽略这些项，相当于引入了不相关的假设。随机森林树之间共享原始数据、交叉验证折之间共享训练记录，都不能没有说明就当作完全独立重复。

\section{范数、正交与投影}
欧氏范数 $\|\mathbf x\|_2=\sqrt{\sum_jx_j^2}$表示当前坐标单位下的长度；内积为0表示正交。单位方向 $\mathbf v$上的投影坐标为 $z=\mathbf v^\mathsf T\mathbf x$，重构投影点为 $z\mathbf v$。方向向量与样本投影值不是同一个对象。

对中心化数据的协方差 $S=X_c^\mathsf TX_c/(n-1)$，$\mathbf v^\mathsf TS\mathbf v$就是该方向投影的样本方差。若标准化把坐标尺改变了，最大方差方向也可能改变；几何公式本身不会替我们决定正确物理尺度。
